\documentclass[11pt,a4paper]{article}
\usepackage[margin=27mm]{geometry}
\usepackage{amsmath,amssymb,amsthm,hyperref}
\newtheorem{theorem}{Theorem}
\newtheorem{lemma}[theorem]{Lemma}
\newcommand{\gen}[1]{\langle #1\rangle}
\newcommand{\F}{\mathbb F}
\newcommand{\Aut}{\operatorname{Aut}}
\title{A chiral coset geometry with disconnected residues}
\author{Human Compression project}
\date{}
\begin{document}\maketitle
\begin{abstract}
We give a rank-four counterexample to the stated residual-connectedness
criterion for chiral geometries constructed from $C^+$-groups. The
specified group is $\operatorname{PSL}_2(13)$, whereas the full
type-preserving automorphism group is
$\operatorname{PSL}_2(13)\times C_2$. The geometry has 112 objects and
4368 chambers. Its full group has two chamber orbits, with adjacent
chambers in different orbits, but a rank-two residue is the disjoint
union of two six-cycles. We prove the full automorphism assertion by
an elementary calculation on thirteen points; no automorphism-group
software or classification is required.
\end{abstract}

\section{The decorated point graph}
All entries in this note are in $\F=\F_{13}$. Let
\[
 K=\{1,3,4,9,10,12\}=(\F^\times)^2,
 \qquad V=(\F^2\setminus\{0\})/K.
\]
Write $[v]$ for the class $Kv$. There are 28 classes, two on each
projective line. Define a simple undirected graph $X$ on $V$ by
\[
                  [v]\sim[w]\quad\Longleftrightarrow\quad
                  \det(v,w)\in K.                         \tag{1}
\]
The condition is independent of the representatives and symmetric
because $-1\in K$. Points on the same projective line are not adjacent.
Put $x=[(1,0)^T]$, $y=[(0,1)^T]$, and $z=[(1,1)^T]$.

The natural $\operatorname{SL}_2(13)$ action preserves (1). Its kernel
is $\{\pm I\}$: a matrix fixing $x,y,z$ must be scalar, and a
determinant-one scalar has square one. Hence its image
$G=\operatorname{PSL}_2(13)$ has order $13(13^2-1)/2=1092$.
It is transitive on vertices and on ordered edges. Indeed a nonzero
vector can be completed to a determinant-one basis; for an adjacent
ordered pair one can first rescale one vector by a member of $K$ to
obtain the required determinant. Every vertex has thirteen neighbours:
those of $x$ have unique representatives $(a,1)^T$, $a\in\F$.
The graph induced on this neighbourhood is
\[
          a\sim b\quad\Longleftrightarrow\quad a-b\in K,
                                                        \tag{2}
\]
the Paley graph on thirteen points. In particular every edge of $X$
has six common neighbours.

Multiplication of every vector by 2 induces an automorphism $s$ of
$X$, since $4\in K$. It is an involution and commutes with $G$.
It does not belong to $G$. Otherwise a determinant-one matrix inducing
this action would be $2kI$ for some $k\in K$, forcing
$2k=\pm1$, whereas $\pm7\notin K$. Thus
\[
             A=\gen{G,s}\cong G\times C_2,
             \qquad |A|=2184.                            \tag{3}
\]
Every element of $A$ is induced by an invertible linear matrix, taken
modulo the scalar subgroup $K$.

\begin{lemma}\label{full-aut}
The full automorphism group of $X$ is $A$.
\end{lemma}
\begin{proof}
The other decorated point on the line of $[v]$ is $[2v]$. For each
different line, exactly one of its two decorations is adjacent to
$[v]$, and the other is adjacent to $[2v]$. These two vertices
therefore have no common neighbour. Two vertices on distinct lines
do have a common neighbour: solve the independent linear equations
$\det(v,w)=\det(v',w)=1$. Thus $[2v]$ is the unique other vertex
with no common neighbour with $[v]$, and every automorphism preserves
this antipodal pairing.

An automorphism fixing $x$ and its thirteen neighbours fixes their
antipodes and the antipode of $x$, hence every vertex. Consequently
the stabilizer of $x$ embeds in the automorphism group of (2).
We bound the latter group directly. The neighbours of 0 in (2)
induce the six-cycle, in the following cyclic order:
\[
                    (1,4,3,12,9,10).                     \tag{4}
\]
Number its positions $0,\ldots,5$. The six remaining nonzero
vertices have the following neighbour sets on this cycle:
\[
\begin{array}{c|rrrrrr}
\text{vertex}&2&5&6&7&8&11\\\hline
\text{positions}&023&014&245&125&134&035.
\end{array}                                               \tag{5}
\]
The six triples in (5) are the translates of $\{0,2,3\}$ in
$\mathbb Z/6\mathbb Z$. They are distinct, so an automorphism
fixing 0 is determined by its action on (4). This action belongs
to the dihedral group of the cycle. Every rotation preserves
the collection (5), while the reflection $i\mapsto-i$ sends
$\{0,2,3\}$ to $\{0,3,4\}$, which is absent. No reflection
preserves the collection. The stabilizer of 0 in the automorphism
group of (2) therefore has order at most six, and that group
has order at most $13\cdot6=78$.

It follows that $|\Aut(X)|\le28\cdot78=2184$. The subgroup (3)
attains the bound, proving equality.
\end{proof}

There are $28\cdot13\cdot6=2184$ ordered triangles in $X$.
An element of $A$ fixing one of them fixes three different
projective lines and is scalar; fixing the decorated points then
makes it the identity modulo $K$. Thus $A$ acts freely and
transitively on the ordered triangles. In particular each triangle
has the same number of common neighbours as $x,y,z$.
A common neighbour of $x,y$ is represented by $(a,1)^T$ with
$a\in K$. It is adjacent to $z$ precisely when $1-a\in K$.
Directly from the six displayed elements of $K$, the solutions
are $a=4,10$. Every triangle has exactly two common neighbours.

\section{The geometry and its full chirality}
Let $I=\{0,1,2,3\}$. Take four tagged copies of $V$, and declare
objects of different types incident precisely when their underlying
vertices are adjacent in $X$. Equality is the only incidence within
one type. Denote this incidence system by $\Gamma$.
Every incident pair extends to a triangle; every triangle has exactly
two common neighbours. Hence every flag extends to a chamber and
every panel has exactly two completions. Thus $\Gamma$ is a thin
geometry. Its objects number $4\cdot28=112$, and its chambers
number $2184\cdot2=4368$.

\begin{lemma}\label{geometry-aut}
The full type-preserving automorphism group of $\Gamma$ is the
diagonal action of $A$. It has exactly two chamber orbits, and
adjacent chambers belong to different orbits.
\end{lemma}
\begin{proof}
First, distinct vertices of $X$ have distinct neighbourhoods.
Antipodes have disjoint neighbourhoods. If $u,v$ are on different
lines, one of the two decorations on the line of $u$ is adjacent
to $v$, while neither is adjacent to $u$; their neighbourhoods
therefore differ.

Write a type-preserving automorphism as permutations $f_i$ on the
four copies of $V$. For distinct $i,j$, incidence gives
$N(f_i(v))=f_j(N(v))$. Given $i,k$, choose $j$ distinct from both.
The preceding equality gives $N(f_i(v))=N(f_k(v))$, so the absence
of equal neighbourhoods implies $f_i=f_k$. All four permutations
are the same automorphism of $X$. Lemma~\ref{full-aut} proves the
first assertion.

This group is free on chambers, since it is free on ordered
triangles. Its order and the chamber count give exactly two
orbits. If an automorphism took one chamber to an adjacent
chamber, it would fix the three common objects and hence be
the identity. Adjacent chambers are distinct, so this cannot
happen. This is chirality under the full type-preserving group,
not merely under a proposed subgroup.
\end{proof}

The geometry is not residually connected. Fix the type-1 object
$x$ and the type-2 object $y$. The remaining objects in their
residue are two copies, of types 0 and 3, of
\[
                   \{[(a,1)^T]:a\in K\}.
\]
Their incidence is $a-b\in K$, so this residue is the bipartite
double of the cycle (4). Since that cycle has even length, its
bipartite double is two disjoint six-cycles: one component uses
the even cycle positions in type 0 and the odd positions in
type 3, and the other uses the opposite positions.

\section{The exact \texorpdfstring{$C^+$}{C-plus} coset construction}
It remains to verify the algebraic hypothesis in the source theorem.
In $G=\operatorname{PSL}_2(13)$ use the determinant-one matrices
\[
 \alpha_0=I,\qquad
 \alpha_1=\begin{pmatrix}4&0\\7&10\end{pmatrix},\quad
 \alpha_2=\begin{pmatrix}4&6\\0&10\end{pmatrix},\quad
 \alpha_3=\begin{pmatrix}3&0\\0&9\end{pmatrix}.             \tag{6}
\]
Choose the following base objects, in type order:
\[
       v_0=[(4,1)^T],\qquad v_1=x,
       \qquad v_2=y,\qquad v_3=z.                        \tag{7}
\]
They form a chamber, as their six determinants belong to $K$.
For $S\subseteq I$ put
$P_S=\gen{\alpha_i^{-1}\alpha_j:i,j\in S}$.
The six differences for $i<j$ are as follows; the last column
gives their two fixed base objects.
\[
\begin{array}{c|c|c}
ij&\alpha_i^{-1}\alpha_j&\text{fixed labels}\\\hline
01&\left(\begin{smallmatrix}4&0\\7&10\end{smallmatrix}\right)&2,3\\
02&\left(\begin{smallmatrix}4&6\\0&10\end{smallmatrix}\right)&1,3\\
03&\left(\begin{smallmatrix}3&0\\0&9\end{smallmatrix}\right)&1,2\\
12&\left(\begin{smallmatrix}1&8\\11&11\end{smallmatrix}\right)&0,3\\
13&\left(\begin{smallmatrix}4&0\\5&10\end{smallmatrix}\right)&0,2\\
23&\left(\begin{smallmatrix}4&11\\0&10\end{smallmatrix}\right)&0,1
\end{array}                                               \tag{8}
\]
Each is non-scalar and its cube is $I$ or $-I$, so it has order
three in $G$. Each fixes the two displayed decorated points,
not just their projective lines. No nonidentity projective
transformation fixes three distinct lines.

The stabilizer in $G$ of any decorated point is $C_{13}:C_3$,
of order 39. For $x$ this follows directly from its upper
triangular lifts, whose first diagonal entry belongs to $K$;
the others follow by transitivity. For each $i$, (8) puts
$P_{I\setminus\{i\}}$ in this stabilizer of $v_i$ and supplies
two distinct subgroups of order three, distinguished by their
second fixed lines. A subgroup of $C_{13}:C_3$ containing two
such subgroups must be the whole group: if it met the normal
$C_{13}$ trivially it would embed in $C_3$, and a nontrivial
intersection is the whole $C_{13}$. Thus
\[
             H_i:=P_{I\setminus\{i\}}=G_{v_i},
             \qquad |H_i|=39.                            \tag{9}
\]

The three matrices generate $G$. Indeed $H_1$ contains every
upper unipotent matrix and $H_2$ every lower unipotent matrix.
These generate $\operatorname{SL}_2(13)$ before passing to its
projective quotient. For instance the explicit products
\[
 \alpha_2\alpha_3=\begin{pmatrix}-1&2\\0&-1\end{pmatrix},
 \qquad
 \alpha_1\alpha_3=\begin{pmatrix}-1&0\\8&-1\end{pmatrix}
\]
already give nonzero upper and lower translations, whose powers
give all such translations. Row elimination then gives every
determinant-one matrix.

The intersection property is now explicit. Two distinct groups
in (9) intersect in a group of order three, since $G$ is transitive
on the $28\cdot13$ ordered edges. Their intersection is the
group $P_S$ from (8) on their two complementary labels. Three
distinct groups in (9) intersect trivially. A cyclic $P_S$ with
$|S|=2$ is contained in $H_i$ exactly when $i\notin S$;
otherwise an element of the intersection would fix three base
lines. Two different such cyclic groups likewise intersect
trivially. Together with $P_I=G$ and $P_S=1$ for $|S|\le1$,
these observations prove
\[
                       P_S\cap P_T=P_{S\cap T}
                       \qquad(S,T\subseteq I).           \tag{10}
\]
The generators are independent: omitting $\alpha_i$, $i>0$,
leaves $H_i$ of order 39. Hence (6) is a $C^+$-marking in the
source's exact convention.

Finally the source's right coset $H_i g$ corresponds to
$g^{-1}v_i$. Equality of fibres follows from (9). Coset
incidence is precisely the $G$-orbit of the corresponding
base pair; every base pair in (7) is an edge, and $G$ is
transitive on all ordered edges. Therefore all six coset
incidence relations agree with (1). This proves that $\Gamma$
is exactly Construction 3.2 for (6), rather than a geometry
which merely contains a suitable group action.

\begin{theorem}
The $C^+$-group (6) has a chiral coset geometry that is not
residually connected. Its full type-preserving group is
$\operatorname{PSL}_2(13)\times C_2$, its specified group is
$\operatorname{PSL}_2(13)$, and its rank-two residue at the
type-1 and type-2 objects of (7) is two disjoint six-cycles.
\end{theorem}

\section{Source scope and verification}
The theorem contradicts Theorem 1.1 in the fixed published
edition of Leemans--Tranchida~\cite{source}. Their definition
on printed page 493 uses the full $\Aut_I(\Gamma)$ to define
chirality; Lemma~\ref{geometry-aut} verifies exactly that
definition. Construction 3.2 on page 494 is paid by (6)--(10)
and the coset identification. The geometry is thin and every
flag extends, so neither of these hypotheses is missing.

The issue is the passage on printed page 496 from proper-flag
transitivity under the full group to transitivity under the
specified $G$. Here the full group acts freely and transitively
on flags of any three fixed types, while $G$ has two orbits on
those flags. The internal rank-two coset model thus covers
only one of the two actual residue cycles. This counterexample
does not refute the separately valid rank-three conclusion in
Proposition 4.11, nor a criterion which explicitly requires
the specified group to have the two alternating chamber orbits.
It is a correction of a published theorem, not a claim to
have solved a named open conjecture.

The decorated-vector vocabulary comes from the accompanying
recovery of the tetrahedral construction of Catalano, Fernandes,
Hubard and Leemans~\cite{tetrahedral}. That source uses point
and tetrahedron types. The present construction instead takes
four point types. All group and incidence arguments needed
for this example are given here. We make no claim of priority
for the graph (1), the Paley graph, or their standard matrix
actions.

The accompanying short program independently enumerates the
28 vertices, all triangles and four-cliques, the generated
1092-element matrix group and all 256 intersections in (10).
It checks the full ordered-triangle action, both chamber
orbits, the six local attachment triples and the disconnected
actual residue. These are finite checks of the displayed
example; the full automorphism upper bound is the written
argument of Lemma~\ref{full-aut}.

\begin{thebibliography}{9}
\bibitem{source} D.~Leemans and P.~Tranchida,
\emph{On residual connectedness in chiral geometries},
Algebraic Combinatorics 4 (2021), 491--499,
\href{https://doi.org/10.5802/alco.162}{doi:10.5802/alco.162}.
\bibitem{tetrahedral} D.~Catalano, M.~E.~Fernandes, I.~Hubard and
D.~Leemans, \emph{Hypertopes with tetrahedral diagram},
Electronic Journal of Combinatorics 25(3) (2018), P3.22.
\end{thebibliography}
\end{document}
